\chapter{Formal Specification of Rego}
\label{ch:form}

This chapter provides formal specification of the Rego policy language used in OPA.
These definitions are derived from the official OPA reference \cite{opa:web} and the OPA source code \cite{opa:git}.
All examples of Rego policies are valid code, interpretable online at the official Rego Playground: \url{https://play.openpolicyagent.org}.

%%%%%%%%%%%%%%%%%%%%%%%%%%%%%%%%%%%%%%%%%%%%%%%%%%%%%%%%%%%%%%%
% 80%
% variables, atomic values, composite values
\section{Terms and Values}
\label{s:terms}

%%===

This section establishes the vocabulary of values, variables, and terms that \rego policies operate over.
We define a substitution $\sigma$ mapping terms to their values; the rest of the chapter is built on this foundation.

%%===

\subsection{Variables and Substitution}
\label{ss:vars}

%%---

A \emph{variable} is an identifier matching the regular expression $[\_a-zA-Z][\_a-zA-Z0-9]^*$.
The symbol $\_$ alone is reserved as a placeholder for an \emph{unnamed variable}: a variable that is intentionally left unbound.
We denote the set of all variable names with $\Vars$.

Variables in a \rego policy are either \emph{global} or \emph{local}.
Global variables are accessible throughout the whole policy and their values compose the response document.
Local variables are scoped to the rule in which they are declared.
We write $\Globals \subset \Vars$ for the set of global variable names and $\Locals \subset \Vars$ for the set of local variable names, with $\Globals \cap \Locals = \emptyset$.
The internal OPA compiler renames all local variables to unique names prior to evaluation, so name shadowing does not arise in the formalization.

We define the set of \rego \emph{terms} $\Terms$ to encompass literals, variables, references, and function calls.
The central semantic object of this chapter is the \emph{substitution}
$$\sigma : \Terms \to \Vals$$
mapping each term to its value.
The codomain $\Vals$ includes the special element $\Undef$, whose properties are defined in Section~\ref{ss:ref}.
We call a term $t$ \emph{ground} (with respect to $\sigma$) if every variable occurring in $t$ is bound under $\sigma$; we call $t$ \emph{defined} if $\sigma(t) \neq \Undef$.
Every defined term is ground, but a ground term may still be undefined — for example, when a reference accesses a missing key in a collection.

The restriction $\sigma_V = \sigma|_{\Globals}$ to global variable names is precisely the response document returned by OPA to the querying service.

%%===

\subsection{Values}
\label{ss:vals}

%%---

The set $\Vals$ of \rego values is defined by the grammar:
\begin{align*}
v \;::=\;\;
  & s \in \Str
    & \text{(string literal)} \\
  \mid\;\; & n \in \Int
    & \text{(number)} \\
  \mid\;\; & b \in \Bool
    & \text{(boolean)} \\
  \mid\;\; & \Null
    & \text{(null)} \\
  \mid\;\; & [v_0, \ldots, v_{n-1}]
    & \text{(array)} \\
  \mid\;\; & \{v_{k_1}{:}v_1, \ldots, v_{k_n}{:}v_n\}
    & \text{(object)} \\
  \mid\;\; & \{v_1, \ldots, v_n\}
    & \text{(set)} \\
  \mid\;\; & \Undef
    & \text{(undefined)}
\end{align*}

Values of types $\Str, \Int, \Bool$, and $\Null$ are called \emph{atomic}; arrays, objects, and sets are called \emph{composite}.
$\Undef$ is a distinguished element that does not belong to any of the above domains; its formal properties are given in Section~\ref{ss:ref}.
The empty array is written $[\,]$, the empty object as $\{\}$, and the empty set as $\texttt{set()}$.
All composite types in \rego are \emph{heterogeneous}: their elements and field values may be of different types.

A \emph{literal} is a term whose value is determined syntactically.
The substitution maps every literal to itself: $\sigma(\ell) = \ell$.
A composite literal is evaluated element-wise; if any constituent term evaluates to $\Undef$, the entire composite literal evaluates to $\Undef$.

A variable $x \in \Vars$ is \emph{bound} under $\sigma$ if $\sigma(x) \neq \Undef$ and \emph{unbound} otherwise.
Assignment of a value to a variable is written $x \mathrel{:=} t$ and yields $\sigma(x) = \sigma(t)$:
\[
\textbf{(Assign)}\quad
\inferrule{
  x \mathrel{:=} t
}{
  \sigma(x) = \sigma(t)
}
\]
\rego provides two syntactic forms for assignment: the assign operator \texttt{:=} and the unification operator \texttt{=}\footnote{
  The unification operator \texttt{=} only functions as an assignment for an undefined left-hand side and it behaves as an equality test otherwise.
}.
This work uses the assign form throughout, in line with the OPA reference~\cite{opa:web:ass}.

%%===

\subsection{References}
\label{ss:ref}

%%---

A \emph{reference} is a term of the form $t[v]$, denoting access into the composite value of $t$ at key or index $v$.
References may be chained: $t[v_1][v_2]\cdots[v_n]$ accesses a nested value.
As syntactic sugar, $t\texttt{.}k$ abbreviates $t[\texttt{"k"}]$ when $k$ is a string literal identifier, so \texttt{x.a.b.c} is identical to \texttt{x["a"]["b"]["c"]}.

The value of a reference is defined recursively:
\[
\sigma(t[v]) \;=\;
\begin{cases}
  \Undef & \text{if } \sigma(t) = \Undef \text{ or } \sigma(v) = \Undef, \\
  \sigma(t)\bigl(\sigma(v)\bigr) & \text{if } \sigma(v) \in \mathrm{dom}\bigl(\sigma(t)\bigr), \\
  \Undef & \text{otherwise,}
\end{cases}
\]
where $\sigma(t)(k)$ denotes the lookup of key $k$ in the value $\sigma(t)$ — array indexing for arrays and field access for objects.
$\mathrm{dom}(v)$ then represents the set of all keys defined in a composite value $v$.
A chained reference evaluates left-to-right; if any intermediate step yields $\Undef$, the entire reference is $\Undef$.

Two special roots exist in every evaluation context.
The identifier \texttt{input} refers to the input document supplied by the querying service.
The identifier \texttt{data} is the root of a hierarchical document that combines the service-provided data (for example, a role-to-permission mapping) with the namespaces of all loaded \rego packages.
For a global variable $x$ defined in package $\pkg$, the reference \texttt{data.$\pkg$.x} and the bare name $x$ denote the same value within the same package: $\sigma(\texttt{data.\textit{$\pkg$}.x}) = \sigma(x)$.

The element $\Undef$ has the following properties that distinguish it from all concrete values:
\begin{itemize}
  \item Equality between two undefined terms is itself undefined: $\Undef$ does not compare equal to $\Undef$.
  \item Any expression containing a subterm that evaluates to $\Undef$ evaluates to $\Undef$.
  \item $\Undef$ is distinct from $\Null$: $\Null$ is a concrete JSON value used to explicitly represent the absence of information (for example, passed as a function argument to signal that an optional parameter is not provided), whereas $\Undef$ means the term is entirely absent from the evaluation result.
  \item undefined global variables do not appear in the output response document.
\end{itemize}

%%===
%% TODO: maybe change set type to contain values, instead of types
\subsection{Type System}
\label{ss:types}

%%---

Each value in $\Vals$ has a corresponding \emph{type}.
Since \rego is dynamically typed, the type of every term is determined at runtime from the concrete value produced by $\sigma$.
The set of \rego types $\Types$ is defined by the grammar:
\begin{align*}
\ty \;::=\; & \Str \mid \Int \mid \Bool \mid \Null \mid \Undef \\
   \mid\;   & \Array{\ty_0,\ldots,\ty_{n-1}} \\
   \mid\;   & \RegoSet{\ty_1,\ldots,\ty_m} \\
   \mid\;   & \Obj{v_1{:}\ty_1,\ldots,v_n{:}\ty_n}
\end{align*}

The intended meaning of each construct:
\begin{itemize}
  \item Atomic types $\Str, \Int, \Bool, \Null$ classify the corresponding atomic values;
        $\Undef$ classifies the undefined element.
  \item $\Array{\ty_0,\ldots,\ty_{n-1}}$ is an \emph{array type} of length $n$,
        where $\ty_i = \typof(v_i)$ is the type of the element at index $i$.
  \item $\RegoSet{\ty_1,\ldots,\ty_m}$ is a \emph{set type} whose element-type list
        is itself unordered and duplicate-free: $\{\ty_1,\ldots,\ty_m\}$ is the set of
        distinct types that appear among the elements, so $\typof(\{v_1,\ldots,v_k\}) =
        \RegoSet{\typof(v_1),\ldots,\typof(v_k)}$ with duplicates removed.
  \item $\Obj{v_1{:}\ty_1,\ldots,v_n{:}\ty_n}$ is an \emph{object type}
        with keys $v_i \in \Vals$ typed $\ty_i$.
\end{itemize}

We write $\typof : \Terms \to \Types$ for the function mapping each term to its type,
with $\typof(t) = \Undef$ whenever $\sigma(t) = \Undef$.
The analysis of how $\typof$ is approximated from a policy without a concrete evaluation
context is the subject of Chapter~\ref{ch:typ}.

%%%%%%%%%%%%%%%%%%%%%%%%%%%%%%%%%%%%%%%%%%%%%%%%%%%%%%%%%%%%%%%
\section{Rules}
%%=============================================================

Every Rego policy is comprised of a set of rules $\Rules$.
Since it is a declarative language, we can truly treat specified rules as a set without a strict order.
Every rule $\Rule$ consists of a \emph{rule head} ($h$) and a list of \emph{rule bodies} ($B$), where the head specifies which term (non-literal) is set, and the body specifies a value to which it is set, and conditions (expressions) which need to be met for this assignment to be valid.

%%=============================================================
\subsection{Rule Head and Body}
%%-------------------------------------------------------------

Head $h$ is defined as a tuple $(h_n, h_k, h_a)$, where $h_n \in \Terms$ denotes its name, $h_k \in (\Terms \cup \Undef)$ denotes its \emph{key}, and $h_a = (a_1,\dots,a_n)$ is a (possibly empty) list of arguments for $n \in \Nat$.
We will often abuse the notation, and call variables heads.
For a rule with head $h = (h_n, h_k, h_a)$ we define a set of rules
We denote a rule with head $h$ as \emph{rule of $h$} and declare a set of all rules of $h$ as $P_h$.

We define body $B$ as a pair $(B_v,B_E)$, where $B_v \in \Terms$ is its value, and $B_E = (e_1,\dots,e_n)$ for $1 \leq i \leq n$ and $e_i$ being an expression.
When all expressions of rule body $B$ hold, we say that it is \emph{satisfied}, denoted as $\Holds(B)$.
Formally said, for a rule body $B$: 
$$\Holds(B) \iff \forall_{e \in B_E}: \Holds(e)$$.
Rules are defined as pairs $(h,B)$.

Every rule is written as $h$ $B_1$ \texttt{else} $\dots$ \texttt{else} $B_n$, with $h$ being the head and $B_i$ being the bodies, for $1 \leq i \leq n$, formally noted as a tuple $(h,B)$, where $B$ is the list of bodies $(B_1,\dots,B_n)$.
A body is written as a value possibly followed by either as a single expression $e$, or as $\{ e_1 \dots e_n \}$, where expression $e_i$ and $e_{i+1}$ are separated either by a semicolon (\texttt{;}) or a newline character, for $1 \leq i \le n$.
If the value of body $B$ is missing, it is implicitly \true.
If the body does not contain any expressions, it is automatically evaluated as valid.
In adition, the rule body $B_i$ is only evaluated if none of its \emph{predecessors} was satisfied, i.e., iff $\neg\exists_{j \in (0,i)}: $

\begin{figure}
  \label{f:e:body}
  \caption{Example of Rego policy with two rules containing different types of bodies.}
\begin{lstlisting}[language=Rego]
policy example

x if 1 < 2
y := 1 if {
  z := true
  x == z
} else := 2
\end{lstlisting}
\end{figure}

\begin{example}
  For example, consider the Rego policy $\phi$ depicted in figure \ref{f:e:body}.
  $phi$ contains two rules, one being the rule of $x$ and the other being the rule of $y$, meaning $|P_x| = |P_y| = 1$.
  The rule of $x$ contains a single body, which assigns it \true, iff the expression $1 < 2$ holds, which it does, since it is neither $\Undef$, nor \false.
  Thus $\sigma(x) = \true$.

  The rule of $y$ contains two bodies, where the first one assigns $1$ to the head $y$ if its two expressions hold: 1) declaration of local variable $z$, which is assigned to be \true (declaration always holds), and 2) comparison of equality between $x$ and $z$, where $\sigma(x) = \sigma(z) = \true$ and thus it also holds.
  The second body, which would unconditionally assign the value of its only body ($2$) to the head, is not evaluated, since the conditions of its predecessor were satisfied.
\end{example}

%%=============================================================
\subsection{Types of Rules}
%%-------------------------------------------------------------

Rego offers multiple different types of rules:
%%-------------------------------------------------------------
\subsubsection{Complete definitions.}
\emph{Complete definitions} (or \emph{complete rules}) define the whole value of rule head.
The rules in the prior example given in figure \ref{f:e:body} were all complete, since they assigned the values of their bodies directly into the head.
For such rules, we infer the value of the head with:
\[
\textbf{(CompRuleVal)}\quad
\inferrule{
  (h,B) \in \Rules
}{
  \Holds(B) \implies \sigma(h) = B_v
}
\]
A \emph{conflict} happens, when two different rules try to assign a different value to the same head, yielding the whole policy \emph{invalid}.

%%-------------------------------------------------------------
\subsubsection{Incremental definitions.}
% Sets
% Objects (conflicts)
\emph{Incremental definitions}, (or \emph{incremental rules}) define different parts (\emph{keys}) of the same head for each rule.
Going back to the definition of rule head $h = (h_n,h_k,h_a)$, for incremental rules, $h_k \in \Terms$ and $h_a = ()$, i.e., its key is defined and its list of arguments is empty.

Incremental rules are applicable for sets with special \texttt{contains} keyword (omitting rule body values) followed by the key and for objects with the head beign a reference, where only the last term is considered a key.
They are not applicable to arrays, as they need to be defined only with complete rules.
With incremental rules, a conflict can occur when two rules assign different values to the same key, which can effectively only happen for objects.

For terms $t,k \in \Terms$, providing an incremental rule for value $t$ at key $k$ is possible, iff $t$ is not defined completely by any other rule.
In this case, we not only say that $t$ is defined incrementally, but also that the value $t[k]$ is defined completely.

For a policy $\psi$, we note the set of all heads defined incrementaly as $\Inc$ and the set of all heads with a complete defintion as $\Comp$.
Since single head $h$ can be defined either completely or incrementally, we can say that $(h \in \Comp) \iff (h \notin \Inc)$.

\begin{figure}
  \label{f:e:inc}
  \caption{Example Rego policy containing incremental rules.}
\begin{lstlisting}[language=Rego,escapechar=|]
package example.incremental

x contains 1 |\label{inc:x1}|
x contains 2 if true |\label{inc:x2}|

y.a.b.c := false |\label{inc:y}|
# y := { "d": false } |\label{inc:y_}|

z[v] := true if { 
  some v in x 
  v > 1
}
\end{lstlisting}
\end{figure}

\begin{example}
  Consider the policy $\phi$ given in figure \ref{f:e:inc}.
  $\phi$ contains two incremental rules regarding $x$ on lines \ref{inc:x1} and \ref{inc:x2}, with keys $1$ and $2$ respectively.
  From these rules, we can infer $x$ to be a set, containing both keys, since the first one is contained automatically, since the body of this rule is empty, and for the second one, all expressions from its body hold.

  Line \ref{inc:y} contains a rule, which completely defines \texttt{y.a.b.c} to be \false, whilst incrementally defining \texttt{y.a.b} on key \texttt{c}.
  This means that also \texttt{y} and \texttt{y.a} are defined incrementally, and introducing their complete definition, e.g., uncommenting line \ref{inc:y_} would cause a collision.

  The last rule incrementally defines \texttt{z} on a variable key \texttt{v}.
  In the body of this rule, we see a two expressions: 1) a declaration of an existentially quantified variable $v$ over fields of $x$, and 2) a constraint on this variable.
  Informally speaking, this rule says: \textit{\texttt{z} contains any element of \texttt{x} which is greater than 1}.

  Evaluating all rules gives the value mapping $\sigma(x) = \{1,2\}$, $\sigma(y) := \{"a": \{"b": \{"c": \false\}\}\}$, and $\sigma(z) = \{1: \true\}$.
  Thus, the response document would contain three fields representing these variables.
\end{example}

%%-------------------------------------------------------------
% TODO
\subsubsection{Function Definitions.}
Rules can serve to define custom functions.
In this case, the argument list $h_a = (a_1,\dots,a_n)$ of given head is non-empty, for $n \in \Nat, 0 \leq i \leq n$ and $a_i \in \Terms$.
Functions then behave the same as complete definitions, with each argument $a_i$ being a local variable for the given rule, which is set at runtime with each function call.
These user-defined functions allow % TODO TODO TODO
User-defined Rego functions allow for declaration with concrete arguments.

%%-------------------------------------------------------------
% FIXME: default values for functions
\subsubsection{Default Rules.}
Default rule sets a fallback value $v$ for head $h$, where $\Rules_h \in \Comp$, meaning $\sigma(h) = v$, iff $\forall_{(h,B)}: \neg\exists_{b \in B}: \Holds(b)$.
We define a new function $\Def: \Terms \to \Vals$ mapping every head $h$ to its \emph{default} value $\Def(h)$.
Default rules are applicable at most once, only for heads with complete definition.
If a given head $h \in \Comp$ does not have a default rule, we set $\Def(h) = \Undef$.

\begin{example}
  
\end{example}
% TODO: example


